\section*{Chapitre \ref{ch:g2:moon-in-the-sky} --- La Lune dans le ciel}

\begin{solution}{exo:g2:moon-in-the-sky:1}
Bien plus petite que la Terre, et bien plus loin que les nuages ---
plus loin que ne vole aucun avion. Nous savons que des humains
peuvent lui rendre visite parce que des astronautes l'ont fait~: le
voyage leur a pris trois jours dans chaque sens.
\end{solution}

\begin{solution}{exo:g2:moon-in-the-sky:2}
Des cratères --- des cicatrices en forme de cuvette faites par des
rochers venus de l'espace qui se sont écrasés sur la \omterm{def:g2:moon-in-the-sky:moon}{Lune}. Elles
restent parce que la \omterm{def:g2:moon-in-the-sky:moon}{Lune} n'a ni \omterm{def:g2:air-around-us:wind}{vent} ni pluie pour les user.
\end{solution}

\begin{solution}{exo:g2:moon-in-the-sky:3}
Une attrapeuse de \omterm{def:g1:light-and-shadows:source}{lumière}. Le clair de lune est de la \omterm{def:g1:light-and-shadows:source}{lumière} du
Soleil qui rebondit sur la roche grise de la \omterm{def:g2:moon-in-the-sky:moon}{Lune}, comme un mur blanc
éclaire une pièce sans être une lampe.
\end{solution}

\begin{solution}{exo:g2:moon-in-the-sky:4}
Toujours celle qui fait face au Soleil --- le côté ensoleillé est
éclairé, le côté opposé vit sa propre nuit.
\end{solution}

\begin{solution}{exo:g2:moon-in-the-sky:5}
Oui --- de nombreux après-midi, une \omterm{def:g2:moon-in-the-sky:moon}{Lune} pâle est suspendue dans le
ciel bleu. La meilleure preuve pour un camarade sceptique~: sortir un
tel après-midi et la lui montrer du doigt.
\end{solution}

\begin{solution}{exo:g2:moon-in-the-sky:6}
Croissant, moitié, presque pleine, \omterm{ex:g2:moon-in-the-sky:shapes}{pleine lune} --- puis on redescend
par presque pleine, moitié, croissant, jusqu'à la \omterm{def:g2:moon-in-the-sky:moon}{Lune} «~nouvelle~»,
cachée. Tout le spectacle dure environ un mois.
\end{solution}

\begin{solution}{exo:g2:moon-in-the-sky:7}
Un cri ne ferait aucun bruit --- le son a besoin d'\omterm{def:g2:air-around-us:air}{air}. Un
cerf-volant ne pourrait pas voler~: pas d'\omterm{def:g2:air-around-us:air}{air} signifie pas de \omterm{def:g2:air-around-us:wind}{vent},
et rien sur quoi le cerf-volant puisse s'appuyer.
\end{solution}

\begin{solution}{exo:g2:moon-in-the-sky:8}
La même heure et le même dessin honnête chaque soir (le même endroit
aide aussi)~: un journal honnête ne change qu'une chose à la fois ---
la nuit.
\end{solution}

\begin{solution}{exo:g2:moon-in-the-sky:9}
Les maisons sont proches~: passer devant elles change la façon dont
Rosa les voit --- elles défilent et disparaissent. La \omterm{def:g2:moon-in-the-sky:moon}{Lune} est si
loin que tout le trajet ne change absolument rien à la façon dont
elle la voit~: elle reste à son épaule et semble suivre. La
championne du lointain «~suit~» toujours.
\end{solution}

\begin{solution}{exo:g2:moon-in-the-sky:10}
La Terre elle-même. La \omterm{def:g1:light-and-shadows:source}{lumière} du Soleil rebondit sur la Terre
brillante jusqu'au côté sombre de la \omterm{def:g2:moon-in-the-sky:moon}{Lune} et l'éclaire faiblement ---
c'est la \omterm{def:g1:light-and-shadows:source}{lumière} cendrée~: notre planète rend à la \omterm{def:g2:moon-in-the-sky:moon}{Lune} la politesse.
\end{solution}
